\documentclass[10pt]{letter}
%\documentclass[11pt,titlepage]{article}


%%%%%%%%%%%%%%%%%%%%%%%%%%%%%%%%%%%%%%%%%%%%%
%%%%%%%%%%%%%%%%%%%%%%%%%%%%%%%%%%%%%%%%%%%%%
\usepackage[english]{babel}
\usepackage{marvosym}
% \usepackage{ethfont}
\usepackage{times}
%\usepackage{hyperref}
\usepackage{paralist}
\usepackage{framed}
\usepackage{xcolor}
\definecolor{lightgray}{gray}{0.95}
\usepackage{soul}
\usepackage[left=2cm,top=2cm,right=2cm,bottom=2cm,nohead,nofoot]{geometry}
%\usepackage[onlymath,swash,mathlf,textlf,minionint]{MinionPro}  % automatically loads MnSymbol, amsmath and textcomp
 \usepackage{amsmath}
 \usepackage{amssymb}
\usepackage{soul}
\usepackage{enumitem}
\usepackage{csquotes}
\usepackage{comment}

\usepackage{graphicx}
% \usepackage{caption}
\usepackage{floatrow}
\floatsetup{style=ruled,capposition=bottom}

\renewcommand{\familydefault}{\sfdefault}


\renewenvironment{leftbar}[1][\hsize]
{%
\def\FrameCommand
{%

    {\hspace{5pt}\color{black}\vrule width 3pt}%
    \hspace{0pt}%must no space.
    \fboxsep=\FrameSep\colorbox{lightgray}%\addcomment{No change done yet.}
}%
\MakeFramed{\hsize#1\advance\hsize-\width\FrameRestore}%
}
{\endMakeFramed}
\setlength{\FrameSep}{8pt}

\newenvironment{reply}
{\noindent \color{black}}
{}

\newenvironment{excerpt}
{
\begin{displayquote}
\color{blue}
}
{
\end{displayquote}
}
%%%%%%%%%%%%%%%%%%%%%%%%%%%%%%%%%%%%%%%%%%%%%
%%%%%%%%%%%%%%%%%%%%%%%%%%%%%%%%%%%%%%%%%%%%%


\usepackage{color}
% \usepackage[pdftex]{graphicx}
\newtheorem{thm}{Theorem}
\newtheorem{theorem}{Theorem}
\newtheorem{proposition}{Proposition}
\newtheorem{corollary}{Corollary}

\newcommand{\change}[1]{\textcolor{blue}{#1}}

\newcommand{\lettersection}[1]{{\Large {\bfseries #1}}}

\signature{Bibo Zhang\\ Ilario Filippini}
\date{}

\renewcommand{\familydefault}{\sfdefault}


\begin{document}
% If you want headings on subsequent pages,
% remove the ``%'' on the next line:
% \pagestyle{headings}

\begin{letter}{Associate Editor\\
Anonymous Reviewers\\
IEEE/ACM Transactions on Networking\\
\vspace{0.2cm}}


\opening{Dear Editor and Reviewers,}
\vspace{0.1cm}


We would like to thank you for all your feedback on our submission. Based on the last reviewers' comments, we have amended our manuscript to reflect the minor revision requests. 

In the following, we provide a point-by-point reply to all the questions, and where and how we have revised the previous submission in direct response to the reviewers’ comments.

Thank you for your time and consideration.

\closing{Sincerely,}

%enclosure listing
\vspace{1cm}

\end{letter}





\newpage

%%%%%%%%%%%%%%%%%%%%%%%%%%%%%%%%%%%%%%%%%%%%%%%%%%%%%%%%%%%%%%%%%%%%%%%%%%%%%%%
%%%%%%%%%%%%%%%%%%%%%%%%%%%%%%%%%%%%%%%%%%%%%%%%%%%%%%%%%%%%%%%%%%%%%%%%%%%%%%%
\vspace{10mm}
\lettersection{Reviewer 1}

\vspace{5mm}
\noindent \textbf{\large Comment \#1}
\begin{leftbar}
About full-duplex in 3GPP, it is still not explicit whether 3GPP already supports it. The authors, instead of pointing to the precise section of the 3GPP document that standardizes it (if any), included a general comment and pointed to two articles. The authors should verify and make clear whether they are dealing with a feature supported by the current standards, and if so, precisely point out 3GPP standard and section(s) of interest. Otherwise, they should state clearly that it is a feature not yet supported by the standards.
\end{leftbar}
\begin{reply}
We would like to thank the reviewer for the feedback and apologize for not having been sufficiently clear. As mentioned in our previous reply, 3GPP has included the support for IAB simultaneous transmission and reception in Release 17 (3GPP TS 38.174 version 17.0.0 Release 17). Full-duplex operations, which were traditionally considered very hard technical challenges, are now more and more present in the standardization process, as confirmed by Network-Controlled Repeater (NCR) devices currently under standardization.

We have added a footnote in Section I Introduction of the new version to make this aspect more precise and explicit.
\end{reply}

\vspace{5mm}
\noindent \textbf{\large Comment \#2}
\begin{leftbar}
Regarding the questions about empty buffers and their impact on RA dynamics (no transmission and starvation cf. Comment \#09 in the response letter), it seems no comment was added about this in the paper. Please clarify this aspect in the manuscript as it was done in the response letter.
\end{leftbar}
\begin{reply}
We apologize for missing this point in the revised paper. We have added a new sentence about empty buffers in Section III.B to better explain the downlink starvation problem.
\end{reply}


%%%%%%%%%%%%%%%%%%%%%%%%%%%%%%%%%%%%%%%%%%%%%%%%%%%%%%%%%%%%%%%%%%%%%%%%%%%%%%%
%%%%%%%%%%%%%%%%%%%%%%%%%%%%%%%%%%%%%%%%%%%%%%%%%%%%%%%%%%%%%%%%%%%%%%%%%%%%%%%
\vspace{10mm}
\lettersection{Reviewer 2}

\vspace{5mm}
\noindent \textbf{\large Comment \#1}
\begin{leftbar}
My comments have been addressed
\end{leftbar}
\begin{reply}
We would like to thank the reviewer for the time and effort she/he dedicated to our work. Thanks to her/his valuable comments, we have significantly improved the quality of our article.
\end{reply}

\end{document}
